{\small
\resizebox{\textwidth}{!}{%
\begin{tabular}{lcccc}
\toprule
\textbf{Mutation} & \textbf{Codon Changes} & \textbf{Score} & \textbf{Class} & \textbf{Fitness} \\
\midrule
N57H       & 1 & 2.5 & Low    & Low \\
M66I       & 1 & 1.5 & Low    & High \\
Q67H       & 1 & 2.5 & Low    & Low \\
K70R       & 1--2 & 1.0--3.0 & Low    & Low \\
K70N       & 1 & 2.5 & Low    & Moderate \\
N74D       & 1 & 2.5 & Low    & Moderate \\
A105T      & 1--2 & 2.5--4.5 & Low--Moderate & Low \\
A105V      & 1--2 & 1.0--3.0 & Low    & Moderate \\
T107A      & 1 & 2.5 & Low    & Low \\
\midrule
Q67H+K70R  & 2 & 3.0 & Low    & High \\
Q67H+N74D  & 2 & 3.0 & Low    & High \\
M66I+A105T & 2 & 3.0 & Low    & High \\
M66I+T107A & 2 & 3.0 & Low    & High \\
A105T+T107A & 2 & 3.0 & Low    & High \\
\bottomrule
\end{tabular}%
}

\vspace{0.2cm}
\par\noindent\footnotesize\textit{Note:} Barrier score $= n_{\text{transversions}} \times 2 \allowbreak + n_{\text{transitions}} \allowbreak \pm \sum \text{property\_penalties}$. A105T and K70R have subtype-specific codon variations that increase barrier score in some genetic backgrounds (e.g., A105T: 4.5 in subtype C; K70R: 3.0 in subtype C via AAG$\rightarrow$AGA). Barrier classification: High ($\geq$7), Moderate (4--6), Low ($<$4).
}
